\begin{figure}[H]
\centering
\begin{tikzpicture}[
  >=Latex,
  every node/.style={font=\small},
  box/.style={
    draw,
    rounded corners=2pt,
    align=center,
    inner sep=5pt,
    minimum height=0.92cm,
    text width=2.55cm
  },
  tinybox/.style={
    draw,
    rounded corners=2pt,
    align=center,
    inner sep=4pt,
    minimum height=0.82cm,
    text width=2.0cm
  },
  flow/.style={->,thick},
  stochastic/.style={->,thick,dashed}
]

% ---------------- Direct sampling ----------------
\begin{scope}[shift={(-6.15,0)}]
\node[font=\normalsize\bfseries] at (4.75,2.05)
  {(a) 直接的参数依赖采样};

\node[tinybox] (x1) at (0,0.55) {$\mathbf{x}$};
\node[box] (enc1) at (3.0,0.55) {
  编码器\\[-1pt]
  $\boldsymbol{\mu}_{\phi}(\mathbf{x}),
   \boldsymbol{\sigma}_{\phi}(\mathbf{x})$
};
\node[box] (draw1) at (6.25,0.55) {
  随机采样\\[-1pt]
  $\mathbf{z}\sim q_{\phi}(\mathbf{z}\mid\mathbf{x})$
};
\node[box] (dec1) at (9.5,0.55) {
  解码器\\[-1pt]
  $p_{\theta}(\mathbf{x}\mid\mathbf{z})$
};

\draw[flow] (x1) -- (enc1);
\draw[stochastic] (enc1) -- (draw1);
\draw[flow] (draw1) -- (dec1);

\node[align=center,text width=8.8cm] at (4.75,-0.95) {
把采样视为黑盒随机操作，会遮蔽通往编码器参数 $\phi$ 的常规路径式反向传播通路。
};
\end{scope}

% ---------------- Reparameterised sampling ----------------
\begin{scope}[shift={(-6.15,-4.55)}]
\node[font=\normalsize\bfseries] at (4.75,2.05)
  {(b) 重参数化采样};

\node[tinybox] (x2) at (0,0.55) {$\mathbf{x}$};
\node[box] (enc2) at (3.0,0.55) {
  编码器\\[-1pt]
  $\boldsymbol{\mu}_{\phi}(\mathbf{x}),
   \boldsymbol{\sigma}_{\phi}(\mathbf{x})$
};
\node[tinybox] (eps) at (3.0,-1.05) {
  $\boldsymbol{\epsilon}\sim
  \mathcal{N}(\mathbf{0},\mathbf{I})$
};
\node[box] (combine) at (6.25,0.55) {
  确定性变换\\[-1pt]
  $\mathbf{z}
  =
  \boldsymbol{\mu}_{\phi}
  +
  \boldsymbol{\sigma}_{\phi}
  \odot\boldsymbol{\epsilon}$
};
\node[box] (dec2) at (9.5,0.55) {
  解码器\\[-1pt]
  $p_{\theta}(\mathbf{x}\mid\mathbf{z})$
};

\draw[flow] (x2) -- (enc2);
\draw[flow] (enc2) -- (combine);
\draw[flow] (eps) -- (combine);
\draw[flow] (combine) -- (dec2);

\node[align=center,text width=8.8cm] at (4.75,-2.15) {
在固定 $\boldsymbol{\epsilon}$ 时，$\mathbf{z}$ 是编码器输出的确定性可微函数，因此显式提供了通往 $\phi$ 的路径式梯度。
};
\end{scope}

\end{tikzpicture}
\caption{\textbf{重参数化技巧把随机性与编码器参数分离。}
直接从 $\mathbf{z}\sim q_{\phi}(\mathbf{z}\mid\mathbf{x})$ 采样，会使穿过采样步骤的常规路径式梯度变得不明显。改为采样与参数无关的 $\boldsymbol{\epsilon}\sim\mathcal{N}(\mathbf{0},\mathbf{I})$，并令 $\mathbf{z}=\boldsymbol{\mu}_{\phi}(\mathbf{x})+\boldsymbol{\sigma}_{\phi}(\mathbf{x})\odot\boldsymbol{\epsilon}$，既保持同一个高斯分布，又使固定噪声样本下的 $\mathbf{z}$ 对编码器输出可微。}
\label{fig:vae-reparameterisation}
\end{figure}